% !TEX root = chapter-standalone.tex

\section{Relations}
\label{sec:connection-relations}

\linkvideo{spring2021-relations:relations:rel-def} % Definition of relation
A basic mathematical notion that underlies the above discussion is that of a \textbf{binary relation}.

\begin{ctdefinition}[Binary relation]
    \label{def:binary-relation}
    A \maindef{binary relation} \relA from a set~\setA to a set~\setB is a subset of the \SY{cartesian product}~$\setA\cartprod \setB$:
    \begin{equation}
        \relA \setsubseteq \setA \cartprod \setB.
    \end{equation}
\end{ctdefinition}

We will often drop the word ``binary'' and simply use the name ``relation''.

We also write
\begin{equation}
    \relA \colon \setA \sto \setB
\end{equation}
to indicate a relation from~\setA to~\setB.
(\setA is the source, and~\setB is the target).

\begin{marginfigure}
    \centering
    \includesag{30_rel_1}
    \caption{Relation of \cref{exa:simple-rel}.}
    \label{fig:example_rel}
\end{marginfigure}

\begin{example}
    \label{exa:simple-rel}
    Let~$\setA = \makeset{ \sbretzel, \sfondue, \schoco}$ and~$\setB = \makeset{ \stea, \smilk, \swater, \swine }$.
    An example of a relation is the subset
    \begin{equation}
        \relA = \makeset{ \tup{\sbretzel, \stea}, \tup{\sfondue, \swater}, \tup{\sfondue, \swine} } \setsubseteq \setA \cartprod \setB.
    \end{equation}
\end{example}

If~\setA and~\setB are finite sets, we can depict a relation~$\relA \colon \setA \sto \setB$ graphically as in \cref{fig:example_rel}.
For each element~$\tup{\ela,\elb} \setin \setA \cartprod \setB$, we draw an arrow from~$\ela$ to~$\elb$ if and only if~$\tup{\ela,\elb} \setin \relA$.

\begin{marginfigure}
    \centering
    \includesag{30_rel_graph}
    \caption{Relations visualized in ``coordinate systems''.}
    \label{fig:example_rel_coord}
\end{marginfigure}

We can also depict this relation graphically as a subset of~$\setA \cartprod \setB$ in a ``coordinate system way'', as in \cref{fig:example_rel_coord}.
The shaded area is the subset~\relA defining the relation.

\begin{remark}[Notation for relations]
    From now on we will also use the following notation, where we write
    \begin{equation}
        \inrel \ela \relA \elb \ \definedas\  \tup{\ela,\elb}\setin \relA
    \end{equation}
    instead of writing~$\tup{\ela,\elb}\setin \relA$.
\end{remark}

\begin{exercise}
    Given an arbitrary set $\setB$, does there always exist a relation $\Emptyset \mto \setB$?
\end{exercise}

\begin{solution}
    Yes.
    Such a relation would be of the form~$\relA \setsubseteq \Emptyset \cartprod \setB = \Emptyset$, where $\setB$ here is an arbitrary set.
    In this situation, $\relA=\Emptyset$ is a relation $\Emptyset \mto \setB$.
\end{solution}

\begin{exercise}
    Given an arbitrary set $\setA$, does there always exist a relation $\setA \mto \Emptyset$?
\end{exercise}

\begin{solution}
    Yes.
    Such a relation would be of the form~$\relA \setsubseteq \setA \cartprod \Emptyset = \Emptyset$, where $\setA$ is an arbitrary set.
    In this situation, $\relA=\Emptyset$ is a relation $\setA \mto \Emptyset$.
\end{solution}

\begin{gradedexercise}[\exname{VisualizeLeqRelation}]
    \label{ex:visualize-leq-relation}
    Let~$\setA = \setB = \makeset{1, 2, 3, 4}$ and consider the relation~$\relA \colon \setA \sto \setB$ defined by
    \begin{equation}
        \relA = \makeset{ \tup{\ela,\elb} \setin \setA \cartprod \setB \mid \ela \leq \elb }.
    \end{equation}
    %
    Visualize the relation~\relA via the method in \cref{fig:example_rel} and \cref{fig:example_rel_coord} each.
\end{gradedexercise}

\solutionof{VisualizeLeqRelation}


\subsection{Properties of relations}
\linkvideo{spring2021-relations:relations:prop-rel} % Properties of relations

% \linkvideo{spring2021-tradeoffs:relations-properties} % Recap of relations
We have seen that relations generalize functions -- every function defines a relation, via its graph, but not every relation comes from a function in this way.
Many notions that we are familiar with for functions also generalize to relations.
Here are a few.

\begin{ctdefinition}[Properties of a relation]
    \label{def:rel_properties}
    \SYNDEF{injective relation}
    \SYNDEF{single-valued}
    \SYNDEF{everywhere-defined}
    \SYNDEF{surjective relation}
    We say that a relation~$\relA \colon \setA \mto \setB$ is:
    \begin{enumerate}
        \item \label{def:injective-relation} \emph{Injective} if
              \begin{equation}
                  \prfsemi{
                      \inrel \ela \relA \elb
                  }{
                      \inrel \elc \relA \elb
                  }{
                      \ela=\elc
                  }
              \end{equation}
        \item \label{def:single-valued-relation} \emph{Single-valued} if
              \begin{equation}
                  \prfsemi{
                      \inrel \ela \relA \elb
                  }{
                      \inrel \ela \relA \eld
                  }{
                      \elb = \eld
                  }
              \end{equation}
        \item \label{def:surjective-relation} \emph{Surjective} if for all~$\elb\setin \setB$ there exists an element~$\ela\setin \setA$ such that~$\inrel{\ela}{\relA}{\elb}$;
        \item \label{def:everywhere-defined-relation} \emph{Everywhere-defined} if for all~$\ela\setin \setA$ there exists an element~$\elb \setin \setB$ such that $\inrel{\ela}{\relA}{\elb}$.
    \end{enumerate}
\end{ctdefinition}

\begin{example}
    The relation depicted in \cref{fig:example_rel} is \SYN{injective relation}{injective} but not \SYN{surjective relation}{surjective}.
    It is not \SY{single-valued}, nor \SY{everywhere-defined}.
\end{example}

